\chapter{Operator Semigroups and Hille--Yosida}

\section{Strongly continuous semigroups and generators}

\begin{notedefinition}
Let $X$ be a Banach space.  A family $(T(t))_{t\ge0}\subseteq\Bop{X}{X}$ is a
\emph{strongly continuous semigroup}, or $C_0$-semigroup, if
\[
  T(0)=I,
  \qquad
  T(t+s)=T(t)T(s),
\]
and
\[
  \lim_{t\downarrow0}T(t)x=x
  \qquad(x\in X).
\]
Its \emph{infinitesimal generator} is
\[
  Ax=\lim_{t\downarrow0}\frac{T(t)x-x}{t},
\]
with domain
\[
  \Dom(A)=\left\{x\in X:
       \lim_{t\downarrow0}\frac{T(t)x-x}{t}
       \text{ exists in }X\right\}.
\]
\end{notedefinition}
\begin{noteproposition}
Let $A$ be the generator of a $C_0$-semigroup $(T(t))$.  Then:
\begin{enumerate}
  \item $A$ is linear and $\Dom(A)$ is dense in $X$;
  \item $T(t)\Dom(A)\subseteq\Dom(A)$ and
        $AT(t)x=T(t)Ax$ for $x\in\Dom(A)$;
  \item for $x\in\Dom(A)$, the orbit is continuously differentiable and
        \[
          \frac{\dd}{\dd t}T(t)x=AT(t)x=T(t)Ax;
        \]
  \item $A$ is closed.
\end{enumerate}
\end{noteproposition}

\begin{proof}
Linearity follows from the limit.  For density, set
\[
  x_h=\frac1h\int_0^hT(s)x\,\dd s.
\]
Strong continuity gives $x_h\to x$, while
\[
  \frac{T(t)x_h-x_h}{t}
  =\frac1h\left(
       \frac1t\int_h^{h+t}T(s)x\,\dd s
       -\frac1t\int_0^tT(s)x\,\dd s
    \right)
  \longrightarrow\frac{T(h)x-x}{h}.
\]
Thus $x_h\in\Dom(A)$.

For $x\in\Dom(A)$,
\[
 \frac{T(h)T(t)x-T(t)x}{h}
 =T(t)\frac{T(h)x-x}{h}\longrightarrow T(t)Ax,
\]
which proves invariance and commutation.  The derivative at arbitrary $t$
follows from the semigroup law; continuity of $s\mapsto T(s)Ax$ makes it a
continuous derivative.

Finally, suppose $x_n\to x$ and $Ax_n\to y$.  The fundamental theorem of
calculus gives
\[
  T(t)x_n-x_n=\int_0^tT(s)Ax_n\,\dd s.
\]
Passing to the limit yields
\[
  T(t)x-x=\int_0^tT(s)y\,\dd s.
\]
Divide by $t$ and let $t\downarrow0$ to obtain $x\in\Dom(A)$ and $Ax=y$.
\end{proof}
\begin{notedefinition}
A $C_0$-semigroup is a \emph{contraction semigroup} if
$\|T(t)\|\le1$ for every $t\ge0$.
\end{notedefinition}

\section{Resolvents of generators}

\begin{notelemma}[Laplace-transform resolvent]
Let $(T(t))$ be a contraction semigroup with generator $A$.  For
$\operatorname{Re}\lambda>0$, define
\[
  R(\lambda,A)x
  =\int_0^\infty e^{-\lambda t}T(t)x\,\dd t.
\]
Then the Bochner integral is well defined,
\[
  R(\lambda,A)=(\lambda I-A)^{-1},
  \qquad
  \|R(\lambda,A)\|\le\frac1{\operatorname{Re}\lambda}.
\]
For real $\lambda>0$,
\[
  \|\lambda R(\lambda,A)\|\le1.
\]
\end{notelemma}

\begin{proof}
The norm estimate follows from
\[
 \int_0^\infty e^{-\operatorname{Re}\lambda t}\|T(t)x\|\,\dd t
 \le\frac{\|x\|}{\operatorname{Re}\lambda}.
\]
For $s>0$, the semigroup law and a change of variables give
\[
 \frac{T(s)-I}{s}R(\lambda,A)x
 =\frac1s\int_0^\infty e^{-\lambda t}
       \bigl(T(t+s)-T(t)\bigr)x\,\dd t
 \longrightarrow \lambda R(\lambda,A)x-x.
\]
Hence $R(\lambda,A)x\in\Dom(A)$ and
$(\lambda I-A)R(\lambda,A)x=x$.  If $x\in\Dom(A)$, integration by parts in
$e^{-\lambda t}T(t)x$ gives
$R(\lambda,A)(\lambda I-A)x=x$.  Thus it is the inverse.
\end{proof}
\begin{notecorollary}
The generator of a contraction semigroup is densely defined and closed,
$(0,\infty)\subseteq\rho(A)$, and
\[
  \|\lambda(\lambda I-A)^{-1}\|\le1
  \qquad(\lambda>0).
\]
\end{notecorollary}

\section{The contraction Hille--Yosida theorem}

\begin{notetheorem}[Hille--Yosida theorem, contraction form]
Let $A\colon\Dom(A)\subseteq X\to X$ be linear.  Then $A$ generates a
contraction $C_0$-semigroup if and only if
\begin{enumerate}
  \item $A$ is densely defined and closed;
  \item $(0,\infty)\subseteq\rho(A)$;
  \item $\|\lambda(\lambda I-A)^{-1}\|\le1$ for every $\lambda>0$.
\end{enumerate}
\end{notetheorem}

\begin{proof}[Sufficiency]
For $\lambda>0$, define the bounded Yosida approximation
\[
  A_\lambda
  =\lambda A(\lambda I-A)^{-1}
  =\lambda^2(\lambda I-A)^{-1}-\lambda I.
\]
Set
\[
  T_\lambda(t)=e^{tA_\lambda}.
\]
Since $\|\lambda(\lambda I-A)^{-1}\|\le1$,
\[
  \|T_\lambda(t)\|
  =\left\|e^{-\lambda t}
     \exp\bigl(t\lambda^2(\lambda I-A)^{-1}\bigr)\right\|
  \le1.
\]
For $x\in\Dom(A)$,
\[
  A_\lambda x=\lambda(\lambda I-A)^{-1}Ax\longrightarrow Ax
\]
as $\lambda\to\infty$.  The resolvent identity shows that the bounded
operators $A_\lambda$ commute.  Duhamel's formula therefore gives, for
$x\in\Dom(A)$,
\[
  \|T_\lambda(t)x-T_\mu(t)x\|
  \le\int_0^t
       \|(A_\lambda-A_\mu)x\|\,\dd s
  =t\|(A_\lambda-A_\mu)x\|,
\]
uniformly for $t$ in compact intervals.  Density of $\Dom(A)$ and the
contraction bounds extend this Cauchy property to every $x\in X$.  Define
\[
  T(t)x=\lim_{\lambda\to\infty}T_\lambda(t)x.
\]
Uniform convergence on compact time intervals gives the semigroup law,
strong continuity, and $\|T(t)\|\le1$.  For $x\in\Dom(A)$,
\[
  T(t)x-x=\lim_{\lambda\to\infty}
          \int_0^tT_\lambda(s)A_\lambda x\,\dd s
        =\int_0^tT(s)Ax\,\dd s,
\]
so the generator extends $A$.  Equality of resolvents for one positive
parameter, or closedness of $A$, shows that the generator is exactly $A$.
Necessity was proved by the Laplace-transform lemma.
\end{proof}
\begin{notecorollary}[Uniqueness and the abstract Cauchy problem]
A generator determines its $C_0$-semigroup uniquely.  If $x_0\in\Dom(A)$,
then
\[
  u(t)=T(t)x_0
\]
is the unique classical solution of
\[
  u'(t)=Au(t),
  \qquad u(0)=x_0.
\]
Moreover, $\Dom(A^2)$ is a core for $A$ and hence is dense in $\Dom(A)$ for
the graph norm.
\end{notecorollary}

\begin{proof}
If $u$ is another classical solution and $T$ is generated by $A$, then for
fixed $t$,
\[
  \frac{\dd}{\dd s}T(t-s)u(s)
  =-AT(t-s)u(s)+T(t-s)Au(s)=0,
\]
so $u(t)=T(t)u(0)$.  For the core statement, apply
$\lambda R(\lambda,A)$ to elements of $\Dom(A)$: these approximants lie in
$\Dom(A^2)$ and converge in the graph norm as $\lambda\to\infty$.
\end{proof}
\begin{notetheorem}[Euler approximation]
If $A$ generates a contraction semigroup, then for every $x\in X$ and $t\ge0$,
\[
  T(t)x
  =\lim_{n\to\infty}
     \left(I-\frac{t}{n}A\right)^{-n}x.
\]
\end{notetheorem}

\begin{proof}
The expression is the exponential approximation associated with the Yosida
operator at $\lambda=n/t$; the convergence follows from the construction in
the Hille--Yosida proof.
\end{proof}
\section{Exponentially bounded semigroups}

\begin{notelemma}
Every $C_0$-semigroup is exponentially bounded: there are $M\ge1$ and
$\omega\in\mathbb R$ such that
\[
  \|T(t)\|\le Me^{\omega t}
  \qquad(t\ge0).
\]
\end{notelemma}

\begin{proof}
Strong continuity and uniform boundedness give
$M_0=\sup_{0\le t\le1}\|T(t)\|<\infty$.  Writing $t=n+s$ with
$n\in\mathbb N_0$ and $0\le s<1$ gives
$\|T(t)\|\le M_0^{n+1}\le M e^{\omega t}$ for suitable $M,\omega$.
\end{proof}
\begin{notelemma}[Higher resolvent powers]
If $\|T(t)\|\le Me^{\omega t}$ and $A$ is its generator, then for
$\operatorname{Re}\lambda>\omega$ and $n\ge1$,
\[
  R(\lambda,A)^n x
  =\frac1{(n-1)!}
    \int_0^\infty t^{n-1}e^{-\lambda t}T(t)x\,\dd t,
\]
and therefore
\[
  \|R(\lambda,A)^n\|
  \le\frac{M}{(\operatorname{Re}\lambda-\omega)^n}.
\]
\end{notelemma}

\begin{proof}
The formula for $n=1$ is the Laplace representation.  Induction, Fubini's
theorem, and the semigroup law reduce the $n$-fold integral over
$t_1+\cdots+t_n=t$ to the simplex volume $t^{n-1}/(n-1)!$.  The estimate is
the Gamma-integral calculation.
\end{proof}
\begin{notetheorem}[Hille--Yosida--Phillips theorem]
Let $A$ be densely defined and closed.  It generates a $C_0$-semigroup
satisfying
\[
  \|T(t)\|\le Me^{\omega t}
\]
if and only if $(\omega,\infty)\subseteq\rho(A)$ and
\[
  \|R(\lambda,A)^n\|
  \le\frac{M}{(\lambda-\omega)^n}
  \qquad(\lambda>\omega,\ n\ge1).
\]
\end{notetheorem}

\begin{proof}[Proof outline]
Necessity is the preceding lemma.  For sufficiency, shift the operator by
$\omega I$ and apply the Yosida construction with the uniform power bounds.
Those bounds replace the contraction estimate and yield uniformly bounded
approximating semigroups after multiplication by $e^{-\omega t}$.  Passing to
the strong limit gives the required semigroup and growth estimate.
\end{proof}
\section{Examples}

\begin{notetheorem}[Translation semigroup]
Let
\[
  X=C_0([0,\infty))
  =\{f\in C([0,\infty)):f(t)\to0\text{ as }t\to\infty\}
\]
with the supremum norm, and define
\[
  (T(t)f)(x)=f(x+t).
\]
Then $(T(t))$ is a contraction $C_0$-semigroup.  Its generator is
\[
  Af=f',
\]
with domain
\[
  \Dom(A)=\{f\in C_0([0,\infty)):
               f\text{ is }C^1\text{ and }f'\in C_0([0,\infty))\}.
\]
\end{notetheorem}

\begin{proof}
Uniform continuity of every $f\in C_0([0,\infty))$ gives strong continuity.
For $f\in\Dom(A)$, the difference quotient converges uniformly to $f'$.  For
the converse, if the uniform difference quotient converges to $g$, then
\[
  f(x+t)-f(x)=\int_0^t g(x+s)\,\dd s,
\]
so $f$ is differentiable with $f'=g\in C_0$.
\end{proof}
\begin{noteproposition}[Positive self-adjoint operators]
Let $H$ be a Hilbert space and let $A$ be densely defined, self-adjoint, and
positive:
\[
  \inner{Ax}{x}\ge0
  \qquad(x\in\Dom(A)).
\]
Then $-A$ generates a contraction $C_0$-semigroup.
\end{noteproposition}

\begin{proof}
For $\lambda>0$,
\[
  \|\lambda x+Ax\|\,\|x\|
  \ge\operatorname{Re}\inner{(\lambda I+A)x}{x}
  \ge\lambda\|x\|^2,
\]
so $\|\lambda(\lambda I+A)^{-1}\|\le1$ whenever the inverse exists.  The
spectral theorem, or the standard range argument for positive self-adjoint
operators, gives $\Ran(\lambda I+A)=H$.  Hille--Yosida applied to $-A$ gives
the claim.  The sign is essential: a positive operator itself generally
produces growth rather than a contraction.
\end{proof}
\section{Dissipative operators and Lumer--Phillips}

\begin{notedefinition}
For $x\in X$, the normalized duality set is
\[
  \mathcal J(x)
  =\{x^*\in X^*:x^*(x)=\|x\|^2,
                    \ \|x^*\|=\|x\|\}.
\]
It is nonempty by Hahn--Banach.  A densely defined operator $A$ is
\emph{accretive} if, for every $x\in\Dom(A)$, there is
$x^*\in\mathcal J(x)$ with
\[
  \operatorname{Re}x^*(Ax)\ge0.
\]
It is \emph{dissipative} if $-A$ is accretive.
\end{notedefinition}

\begin{noteproposition}
An operator $A$ is dissipative if and only if
\[
  \|(\lambda I-A)x\|\ge\lambda\|x\|
  \qquad(x\in\Dom(A),\ \lambda>0).
\]
\end{noteproposition}

\begin{proof}
If $x^*\in\mathcal J(x)$ and
$\operatorname{Re}x^*(Ax)\le0$, then
\[
 \|(\lambda I-A)x\|\,\|x\|
 \ge\operatorname{Re}x^*((\lambda I-A)x)
 \ge\lambda\|x\|^2.
\]
The converse follows by taking a supporting functional for $x$ and
letting $\lambda\downarrow0$ in the norm inequality, or equivalently by the
standard directional-derivative characterization of the norm.
\end{proof}
\begin{notetheorem}[Lumer--Phillips theorem]
Let $A$ be densely defined on a Banach space.  Then $A$ generates a contraction
$C_0$-semigroup if and only if
\begin{enumerate}
  \item $A$ is dissipative;
  \item $\Ran(\lambda_0I-A)=X$ for some $\lambda_0>0$.
\end{enumerate}
In this case the range condition holds for every $\lambda>0$.
\end{notetheorem}

\begin{proof}
A contraction generator is dissipative by differentiating
$\|T(t)x\|\le\|x\|$ at $t=0$ through a supporting functional, and its positive
resolvent is onto.  Conversely, dissipativity makes
$\lambda I-A$ injective with inverse norm at most $1/\lambda$ on its range.
The resolvent identity propagates surjectivity from one positive $\lambda_0$
to all positive $\lambda$ by overlapping Neumann-series intervals.  The
operator is closed, and the contraction Hille--Yosida theorem applies.
\end{proof}
\begin{notecorollary}
Let $A$ be densely defined and closed on a Banach space.  If both $A$ and its
Banach adjoint $A^*$ are dissipative, then $A$ generates a contraction
$C_0$-semigroup.
\end{notecorollary}

\begin{proof}
Dissipativity makes $\Ran(I-A)$ closed.  If it were not dense, Hahn--Banach
would give $0\ne f\in X^*$ annihilating it.  Then $f\in\Dom(A^*)$ and
$(I-A^*)f=0$, contradicting the dissipative inequality for $A^*$.  Hence the
range is dense and closed, so it equals $X$.  Apply Lumer--Phillips.
\end{proof}
\begin{notedefinition}
Let $A$ be closed.  A subspace $D\subseteq\Dom(A)$ is a \emph{core} for $A$ if
the graph of $A|_D$ is dense in $\Graph(A)$, equivalently, if $D$ is dense in
$\Dom(A)$ for the graph norm $\|x\|+\|Ax\|$.
\end{notedefinition}

\begin{notetheorem}[Invariant cores]
Let $A$ generate a $C_0$-semigroup $(T(t))$.  Suppose
$D\subseteq\Dom(A)$ is dense in $X$ and $T(t)D\subseteq D$ for every $t\ge0$.
Then $D$ is a core for $A$.
\end{notetheorem}

\begin{proof}
Let $B$ be the graph closure of $A|_D$.  Invariance and the orbit derivative
show that $B\subseteq A$ and that $B$ generates the restriction of the same
semigroup to its closure.  For sufficiently large $\lambda$, the Laplace
resolvent maps the dense set $D$ into the graph closure of $A|_D$.  Since the
resolvent is bounded and onto $\Dom(A)$ in the graph norm, that closure is all
of $\Dom(A)$.
\end{proof}
\begin{notetheorem}[Gaussian heat semigroup]
On $X=C_0(\mathbb R^d)$, define for $t>0$
\[
  (T(t)f)(x)
  =\frac1{(4\pi t)^{d/2}}
     \int_{\mathbb R^d}
       e^{-|x-y|^2/(4t)}f(y)\,\dd y,
  \qquad T(0)=I.
\]
Then $(T(t))$ is a contraction $C_0$-semigroup.  Its generator is the closure
of
\[
  \Delta\colon C_c^\infty(\mathbb R^d)\to C_0(\mathbb R^d),
\]
and $C_c^\infty(\mathbb R^d)$ is a core.
\end{notetheorem}

\begin{proof}
The Gaussian kernels form an approximate identity and satisfy
$G_t*G_s=G_{t+s}$, which follows from Fourier transform or direct Gaussian
integration.  Positivity and unit mass give the contraction estimate, and
approximate-identity convergence gives strong continuity on $C_0$.  For
$u\in C_c^\infty$, differentiation under the integral or Fourier transform
gives
$\frac{\dd}{\dd t}T(t)u=\Delta T(t)u=T(t)\Delta u$.  The invariant-core
theorem identifies the generator with the closure of $\Delta$ on
$C_c^\infty$.
\end{proof}
\begin{noteremark}
The same formula is not strongly continuous on all of $C_b(\mathbb R^d)$:
functions in $C_b$ need not be uniformly continuous.  The correct standard
spaces are $C_0(\mathbb R^d)$ or the bounded uniformly continuous functions.
\end{noteremark}

\begin{notetheorem}[Poisson semigroup]
Let
\[
  P_t(x)=c_d\frac{t}{(t^2+|x|^2)^{(d+1)/2}},
  \qquad t>0,
\]
where $c_d$ normalizes the integral to one.  Then
\[
  T(t)f=P_t*f
\]
defines a contraction $C_0$-semigroup on $C_0(\mathbb R^d)$.  Its generator
is
\[
  A=-(-\Delta)^{1/2},
\]
so on a suitable core
\[
  A^2=-\Delta.
\]
\end{notetheorem}

\begin{proof}
The Fourier multiplier of $P_t$ is $e^{-t|\xi|}$, so the semigroup property
and contraction estimate follow.  Differentiation at $t=0$ gives the multiplier
$-|\xi|$, which is $-(-\Delta)^{1/2}$.  Squaring gives the multiplier
$|\xi|^2$, corresponding to $-\Delta$.
\end{proof}
